% 中文优化LaTeX模板 - 专门解决中文识别问题
% 智慧水利教材专用

\documentclass[UTF8,12pt,a4paper,oneside]{ctexbook}

% ==== 中文字体配置 ====
% 确保中文字体正确设置
\setCJKmainfont{SimSun}[
    BoldFont=SimHei,
    ItalicFont=KaiTi
]
\setCJKsansfont{SimHei}
\setCJKmonofont{FangSong}

% 英文字体配置
\setmainfont{Times New Roman}
\setsansfont{Arial}
\setmonofont{Consolas}

% ==== 基础包 ====
\usepackage{graphicx}
\usepackage{xcolor}
\usepackage{amsmath}
\usepackage{longtable}
\usepackage{booktabs}
\usepackage{array}

% ==== 页面设置 ====
\usepackage[
    top=2.5cm,
    bottom=2.5cm,
    left=3cm,
    right=2.5cm
]{geometry}

% 行距
\usepackage{setspace}
\onehalfspacing

% ==== 超链接 ====
\usepackage{hyperref}
\hypersetup{
    colorlinks=true,
    linkcolor=blue,
    urlcolor=blue,
    citecolor=blue,
    bookmarksopen=true,
    pdfstartview=FitH,
    unicode=true  % 重要：支持Unicode字符
}

% ==== 代码高亮 ====
\usepackage{listings}
\usepackage{xcolor}

% 定义颜色
\definecolor{codegreen}{rgb}{0,0.6,0}
\definecolor{codegray}{rgb}{0.5,0.5,0.5}
\definecolor{codepurple}{rgb}{0.58,0,0.82}
\definecolor{backcolour}{rgb}{0.95,0.95,0.92}

% 基本代码样式
\lstdefinestyle{mystyle}{
    backgroundcolor=\color{backcolour},   
    commentstyle=\color{codegreen},
    keywordstyle=\color{blue},
    numberstyle=\tiny\color{codegray},
    stringstyle=\color{codepurple},
    basicstyle=\ttfamily\footnotesize,
    breakatwhitespace=false,         
    breaklines=true,                 
    captionpos=b,                    
    keepspaces=true,                 
    numbers=left,                    
    numbersep=5pt,                  
    showspaces=false,                
    showstringspaces=false,
    showtabs=false,                  
    tabsize=2,
    frame=single,
    rulecolor=\color{black},
    extendedchars=false,  % 避免扩展字符问题
    inputencoding=utf8    % 明确指定UTF-8输入
}

\lstset{style=mystyle}

% 特定语言设置
\lstdefinestyle{javascript}{
    style=mystyle,
    language=JavaScript,
    morekeywords={const, let, var, function, async, await}
}

\lstdefinestyle{python}{
    style=mystyle,
    language=Python,
    morekeywords={def, class, import, from, as}
}

\lstdefinestyle{csharp}{
    style=mystyle,
    language=[Sharp]C,
    morekeywords={var, async, await, public, private}
}

% ==== 图片设置 ====
\graphicspath{{images/}}
\setkeys{Gin}{width=0.8\textwidth}

% ==== 标题设置 ====
\ctexset{
    chapter = {
        format = \Large\bfseries\centering\color{blue},
        beforeskip = 20pt,
        afterskip = 30pt
    },
    section = {
        format = \large\bfseries\color{blue},
        beforeskip = 15pt,
        afterskip = 15pt
    },
    subsection = {
        format = \normalsize\bfseries,
        beforeskip = 12pt,
        afterskip = 12pt
    }
}

% ==== 文档信息 ====
\title{\textbf{智慧水利平台架构与开发}}
\author{教材编写组}
\date{\today}

% ==== 文档开始 ====
\begin{document}

% 标题页
\maketitle
\newpage

% 目录
\tableofcontents
\newpage

% 正文内容
\section{第一章
智慧水利概述与平台架构基础}\label{ux7b2cux4e00ux7ae0-ux667aux6167ux6c34ux5229ux6982ux8ff0ux4e0eux5e73ux53f0ux67b6ux6784ux57faux7840}

\subsection{学习目标}\label{ux5b66ux4e60ux76eeux6807}

通过本章学习，学生应能够：

\begin{enumerate}
\def\labelenumi{\arabic{enumi}.}
\tightlist
\item
  理解智慧水利的基本概念、发展背景及其在水利现代化中的重要作用
\item
  掌握智慧水利平台的总体架构体系，包括感知层、传输层、平台层、应用层的功能和关系
\item
  分析国内外智慧水利发展现状与技术趋势，了解智慧水利平台开发的技术路径
\item
  认识软件工程方法在智慧水利平台开发中的应用价值和实践意义
\end{enumerate}

\subsection{引言}\label{ux5f15ux8a00}

智慧水利是新一代信息技术与传统水利工程深度融合的产物，代表着水利行业数字化转型的重要发展方向。随着物联网（Internet
of Things, IoT）、大数据（Big Data）、人工智能（Artificial Intelligence,
AI）、云计算（Cloud
Computing）等技术的快速发展，传统的水利管理模式正在发生深刻变革。智慧水利平台作为这一变革的核心载体，集成了现代信息技术与水利专业知识，为水资源管理、水安全保障、水环境保护等提供了全新的解决方案。

软件工程作为一门应用计算机科学、数学和工程原则来设计、开发、维护和测试软件的学科，为智慧水利平台的开发提供了科学的方法论和技术保障。通过系统化、规范化和量化的软件工程方法，可以确保智慧水利平台的高质量、可靠性和可维护性，有效应对复杂水利系统的开发挑战。

\subsection{本章小节}\label{ux672cux7ae0ux5c0fux8282}

\begin{infobox}[章节结构]

\end{infobox}

\begin{verbatim}
### [第一节 软件概述](section01-01.md)
- 软件的定义与特点
- 软件的分类与发展历程
- 软件危机的概念与影响
- 软件工程的产生背景

### [第二节 智慧水利平台概述](section01-02.md)
- 智慧水利的基本概念
- 智慧水利平台架构体系
- 智慧水利发展现状与趋势
- 智慧水利平台的技术特征

### [第三节 本章小结](section01-03.md)
- 软件工程基础理论总结
- 智慧水利平台核心要点
- 思考题与练习
\end{verbatim}

\subsection{关键概念}\label{ux5173ux952eux6982ux5ff5}

\begin{longtable}[]{@{}
  >{\raggedright\arraybackslash}p{(\linewidth - 4\tabcolsep) * \real{0.3000}}
  >{\raggedright\arraybackslash}p{(\linewidth - 4\tabcolsep) * \real{0.3000}}
  >{\raggedright\arraybackslash}p{(\linewidth - 4\tabcolsep) * \real{0.4000}}@{}}
\toprule\noalign{}
\begin{minipage}[b]{\linewidth}\raggedright
概念
\end{minipage} & \begin{minipage}[b]{\linewidth}\raggedright
定义
\end{minipage} & \begin{minipage}[b]{\linewidth}\raggedright
重要性
\end{minipage} \\
\midrule\noalign{}
\endhead
\bottomrule\noalign{}
\endlastfoot
智慧水利 &
运用物联网、大数据、人工智能等技术，实现水利工程智能化管理的综合体系 &
水利现代化的核心标志 \\
数字孪生 & 通过数字化手段构建物理水利工程的虚拟映射模型 &
提升管理精度和决策水平 \\
物联感知 & 利用传感器网络实时采集水利工程运行数据 &
智慧水利的数据基础 \\
云边协同 & 云计算与边缘计算相结合的混合计算架构 &
平衡性能与成本的关键技术 \\
\end{longtable}

\end{document}